\documentclass[12pt,letterpaper]{article}

\usepackage[margin=1in]{geometry}
\usepackage[T1]{fontenc}
\usepackage[utf8]{inputenc}
\usepackage{lmodern}
\usepackage{microtype}
\usepackage{setspace}
\usepackage{enumitem}
\usepackage{tabularx}
\usepackage{titlesec}
\usepackage{xcolor}
\usepackage[hidelinks]{hyperref}

\onehalfspacing

\definecolor{accent}{HTML}{1A3C6E}

\setlist[itemize]{leftmargin=1.5em,noitemsep,topsep=0pt,parsep=0pt,partopsep=0pt}
\setlist[enumerate]{leftmargin=1.5em,noitemsep,topsep=0pt,parsep=0pt,partopsep=0pt}
\titleformat{\section}{\normalsize\bfseries\color{accent}}{}{0pt}{}
\titlespacing*{\section}{0pt}{0.5ex}{0.05ex}
\setlength{\parindent}{0pt}
\setlength{\parskip}{0.2em}
\setlength{\fboxsep}{3.5pt}
\setlength{\fboxrule}{0.4pt}

\newcommand{\placeholder}[1]{\textcolor{gray}{[\textit{#1}]}}

\begin{document}

\begin{center}
{\large\bfseries The Effect of Feedback Modality on Speed--Accuracy Tradeoff\\[0.05em]in a Visual Target-Selection Task\par}
\vspace{0.05em}
{\small MIE\,286 --- Probability \& Statistics (Winter\,2026) --- Project Proposal\par}
\vspace{0.05em}
{\small\textbf{Team:} \placeholder{Team \#}\hspace{0.6em}\textbf{Members:} \placeholder{Name 1}, \placeholder{Name 2}, \placeholder{Name 3}\hspace{0.6em}February 15, 2026\par}
\end{center}
\vspace{-0.8em}
{\color{accent}\hrule height 0.5pt}
\vspace{0.15em}

\section*{1.\enspace Background and Motivation}

In engineering and UI design, operators receive real-time feedback---production counters, score overlays, timing displays---to guide performance. However, the \emph{modality} of feedback can differentially affect cognitive demands. Wickens' \emph{Multiple Resource Theory} posits that visual and auditory processing use separate attentional pools; when a visual task is paired with visual feedback, interference increases [1,\,2]. Cognitive Load Theory predicts that overloading one channel degrades performance, while distributing across modalities reduces load [3]. In rapid pointing tasks described by Fitts' Law [4,\,5], such interference may shift the \emph{speed--accuracy tradeoff} (SAT)---the relationship where faster responses incur more errors [6]. Prior work shows auditory feedback can reduce time in visual pointing tasks [7] and support motor learning without overloading the visual channel [8], but whether it changes the SAT \emph{structure}---how participants balance speed against accuracy---remains unclear.

\section*{2.\enspace Research Question and Hypotheses}

\textbf{RQ:} Does feedback modality (visual HUD vs.\ auditory tones) influence the SAT in a visual target-selection task?
\begin{itemize}
  \item \textbf{H\textsubscript{1} (Speed):} Mean RT will be lower under auditory feedback, as auditory cues avoid competing for visual resources [1].
  \item \textbf{H\textsubscript{2} (Accuracy):} Auditory feedback will yield higher mean radial error, consistent with a classic SAT where speed gains cost accuracy [6].
  \item \textbf{H\textsubscript{3} (Tradeoff):} The RT--error relationship will differ between conditions, indicating a shifted tradeoff under auditory feedback [6].
\end{itemize}

\section*{3.\enspace Experimental Design and Methods}

\noindent\fcolorbox{accent!50}{accent!3}{%
\begin{minipage}{\dimexpr\linewidth-2\fboxsep-2\fboxrule}
\small\renewcommand{\arraystretch}{1.0}%
\begin{tabularx}{\linewidth}{@{}p{0.48\linewidth}X@{}}
\textbf{IV:} Feedback Type (Visual vs.\ Auditory) & \textbf{Design:} Within-subject; counterbalanced order\\
\textbf{DV\textsubscript{1}:} Response time (ms), onset to click & \textbf{Sample:} $N\!\ge\!24$ adults (18+); $\ge$12 per order\\
\textbf{DV\textsubscript{2}:} Radial error (px), click-to-center dist. & \textbf{Trials:} 2 blocks\,$\times$\,60 (+10 practice)\\
\end{tabularx}
\end{minipage}}

\textbf{Participants.} $\ge$24 adults (18+) with normal/corrected vision and no hearing impairments, recruited via convenience sampling (classmates, friends, family). We record age range, device type (mouse/trackpad), and gaming frequency as covariates. Convenience-sampling limitations will be discussed in the report.

\textbf{Task.} Browser-based target selection on a fixed canvas ($800\!\times\!500$\,px). Each trial shows a circular target (radius\,=\,40\,px) at a random location; participants click as quickly and accurately as possible. After 500\,ms, the next target appears. Coordinates and timestamps are logged.

\textbf{Feedback conditions.}
\emph{Visual:} Border flash (green hit\,/\,red miss, 120\,ms) + persistent HUD (running mean RT, hit rate, mean radial error).
\emph{Auditory:} No HUD (trial counter only); 200\,ms tone after each click---high pitch\,=\,fast hit, medium\,=\,slow hit, low thud\,=\,miss.

\textbf{Procedure ($\boldsymbol{\sim}$15\,min).}
Consent/demographics\,$\to$\,instructions\,$\to$\,10 practice trials\,$\to$\,Block\,1 (60 trials)\,$\to$\,break\,$\to$\,Block\,2 (60 trials). Order counterbalanced by random assignment (V$\to$A vs.\ A$\to$V).

\textbf{Controls.} Target size, canvas, trial count, inter-trial interval, and instructions held constant. Target positions randomized. Within-subject design controls individual differences.

\section*{4.\enspace Analysis Plan (R/RStudio)}
\begin{enumerate}
  \item \textbf{Cleaning:} Exclude RT\,$<$\,200\,ms or $>$\,5\,s; verify trial counts.
  \item \textbf{Aggregation:} Per participant$\times$condition: mean RT, mean radial error, hit rate.
  \item \textbf{Descriptives:} Means, SDs, medians; boxplots of RT and error; RT-vs-error scatterplot by condition.
  \item \textbf{Assumptions:} Shapiro--Wilk on within-subject differences for both DVs.
  \item \textbf{H\textsubscript{1}--H\textsubscript{2}:} Paired $t$-tests (Wilcoxon if non-normal); report $t$, $p$, Cohen's $d$.
  \item \textbf{H\textsubscript{3}:} Pearson $r$ (RT vs.\ error) per condition; \texttt{error\,$\sim$\,RT\,$\times$\,condition}---significant interaction\,$=$\,different SAT slopes.
\end{enumerate}

\textbf{Limitations.} Convenience sampling limits generalizability; device variability and audio environments may add noise (recorded as covariates); fixed difficulty may produce accuracy ceiling effects.

\vspace{0.25em}
{\color{accent}\hrule height 0.4pt}
\vspace{0.08em}

{\scriptsize\setlength{\parskip}{0pt}%
\textbf{References}\par\vspace{-0.2em}
\begin{enumerate}[label={[\arabic*]},leftmargin=1.5em,itemsep=0pt,parsep=0pt,topsep=0pt]
  \item C.\,D.\ Wickens, ``Multiple resources and performance prediction,'' \emph{Theor.\ Issues Ergon.\ Sci.}, 3(2), 159--177, 2002.
  \item C.\,D.\ Wickens, J.\,G.\ Hollands, S.\ Banbury, R.\ Parasuraman, \emph{Eng.\ Psychology \& Human Performance}, 4th\,ed., Routledge, 2015.
  \item J.\ Sweller, ``Cognitive load during problem solving: Effects on learning,'' \emph{Cognitive Science}, 12(2), 257--285, 1988.
  \item P.\,M.\ Fitts, ``The information capacity of the human motor system in controlling the amplitude of movement,'' \emph{J.\ Exp.\ Psychol.}, 47(6), 381--391, 1954.
  \item I.\,S.\ MacKenzie, ``Fitts' law as a research and design tool in human--computer interaction,'' \emph{Hum.--Comput.\ Interact.}, 7(1), 91--139, 1992.
  \item R.\,P.\ Heitz, ``The speed--accuracy tradeoff: History, physiology, methodology, and behavior,'' \emph{Front.\ Neurosci.}, 8, 150, 2014.
  \item M.\ Akamatsu, I.\,S.\ MacKenzie, T.\ Hasbroucq, ``A comparison of tactile, auditory, and visual feedback in a pointing task using a mouse-type device,'' \emph{Ergonomics}, 38(4), 816--827, 1995.
  \item R.\ Sigrist, G.\ Rauter, R.\ Riener, P.\ Wolf, ``Augmented visual, auditory, haptic, and multimodal feedback in motor learning: A review,'' \emph{Psychon.\ Bull.\ \& Rev.}, 20(1), 21--53, 2013.
\end{enumerate}
}

\end{document}
